\scanpage{44}
\begin{center}6 - 24 -\hfill n\textsuperscript{o} 392\end{center}

4)\quad a) Montrer que l'anneau $A=Q[X,Y]/\underline{p}$, où $\underline{p}$ est l'idéal principal engendré par $X^2+Y^2-1$ dans $Q[X,Y]$, est un anneau de Krull non factoriel, (montrer que $X$ est un élément extrémal de $A$ mais que $AX$ n'est pas un idéal maximal de $A$).

b) Montrer que l'anneau $A\otimes_Q Q(i)$ est factoriel, (prouver que cet anneau est isomorphe au quotient de $Q(i)[X,Y]$ par l'idéal principal $(XY-1)$ ; comparer à l'exerc. 3 b)).

5)\quad a) Soient $k$ un anneau factoriel noethérien, $B$ l'anneau de polynômes $k[X_1,\ldots,X_n]$ avec $n\geq 3$ ; soient $g_i$, ($0\leq i\leq r$), des éléments de $k[X_3,\ldots,X_n]$ où $g_0$ est extrémal. On pose $g=X_1X_2-\sum_{i=0}^r g_iX_1^i$, et on considère l'anneau quotient $A=B/gB$. Montrer que $A$ est factoriel. (Soit $S$ la partie multiplicative de $A$ engendrée par $1$ et l'image de $X_1$ dans $A$ ; appliquer à $S^{-1}A$ la prop. 3 du n\textsuperscript{o} 4).

b) Soient $k$ un corps de caractéristique $\ne 2$, $F$ un polynôme homogène du second degré de $k[X_1,\ldots,X_n]$ avec $n\geq 5$, tel que la fonction polynôme correspondante sur $k^n$ soit une forme quadratique non dégénérée. Montrer que l'anneau $k[X_1,\ldots,X_n]/(F)$ est factoriel. (On prouvera d'abord qu'un polynôme homogène $G$ du second degré dans $k[X_1,\ldots,X_n]$ tel que la fonction polynôme correspondante soit une forme quadratique non dégénérée, est extrémal pour $n\geq 3$. Prouver ensuite la proposition lorsque $k$ est algébriquement clos, en utilisant a) ; passer enfin au cas général à l'aide de l'exerc. 3 b)).

c) Si $F=X_1X_2-X_3X_4$, montrer que l'anneau $k[X_1,X_2,X_3,X_4]/(F)$ n'est pas factoriel. (Montrer que les images des $X_i$ dans cet anneau sont des éléments extrémaux).

6)\quad a) Soit $K$ un corps algébriquement clos de caractéristique $\ne 2$. Déterminer les idéaux premiers gradués de l'anneau $K[X,Y,Z]/(X^2+Y^2+Z^2)$.

b) Soient $k$ un corps ordonné, $a$, $b$, $c$ des éléments $>0$ de $k$ et $A$ l'anneau $k[X,Y,Z]/(aX^2+bY^2+cZ^2)$. Montrer que $A$ est un anneau factoriel. (Se ramener au cas où $k$ est un corps ordonné maximal, en utilisant l'exerc. 3 b) ; prouver alors, en utilisant a), que tout idéal premier gradué de hauteur $1$ dans $A$ est principal, et appliquer l'exerc. 3 a)).

\margin{43}

